Three findings qualify conditional forecasts of two-sided absolute holding-period position peaks: FX gains depend on the historical reference and period. Against matched frozen $F_{2880}$, QR lowers loss by \FXMatchedGain{}\% in the inspected period and \RWritingMatchedGain{}\% in the 2015 block with no prior evaluation role found in the repository audit, but earlier loss rises by \RWritingOriginalExcess{}\% against the original primary's rolling $F_{720}$, so that primary's gain direction does not recur; the Treasury $W=52$ analogue raises loss by \RWritingTreasuryExcess{}\%. Across \PanelSeriesCount{} additional weekly series, expanding the ridge menu reduces median excess loss from \PanelOriginalMedian{}\% to \PanelExpandedMedian{}\%, yet \PanelExpandedLosses{} of \PanelCellCount{} cells still lose. On short weekly samples, validation selection, including a conservative one-standard-error rule, does not reliably protect the historical benchmark. Finally, the original FX primary lowers loss by \FXPrimaryGain{}\% while breaches rise from \FXReferenceBreaches{} to \FXPrimaryBreaches{}; the \FXLegacyGain{}\% 24-step legacy gain reverses above an underprediction/overprediction weight ratio of \FXLegacyCrossing{}, and both primary Murphy curves cross. Inventory-limit comparisons therefore require error amounts and breach counts alongside pinball loss. The evidence remains retrospective: prior FX inspection, downloaded-vintage Treasury aggregates and shared panel reporting limit generalization.
